% !TeX root = eleve.tex

% ======== Barème ========
\begin{center}
  \badge{primary}{Ex.1: 4pts} \quad
  \badge{primary}{Ex.2: 4pts} \quad
  \badge{primary}{Ex.3: 8pts} \quad
  \badge{primary}{Ex.4: 4pts} \quad
  \badge{accent}{Bonus: +2pts}
\end{center}

\vspace{0.3cm}

% ======== Consignes ========
\begin{consignebox}
  \faIcon{info-circle} \textbf{CONSIGNES IMPORTANTES}
  \begin{itemize}[leftmargin=*,noitemsep,topsep=5pt]
    \item[\faIcon{check}] Rédigez clairement et justifiez tous vos calculs
    \item[\faIcon{check}] Respectez les priorités opératoires
    \item[\faIcon{check}] Les exercices peuvent être traités dans n'importe quel ordre
    \item[\faIcon{times}] Calculatrice non autorisée
  \end{itemize}
\end{consignebox}

\vspace{0.5cm}

% ======== QCM ========
\begin{exercisebox}{\faIcon{list-ul} QCM : Nombres relatifs \points{4}}
  Pour chaque question, écris la lettre de la seule réponse correcte dans la colonne de droite :
  \vspace{0.3cm}

  % Colonnes: question large, choix centrés, colonne Corrigé masquée en élève
  \newcolumntype{C}[1]{>{\centering\arraybackslash}L{#1}}
  {\small\setlength{\tabcolsep}{4pt}\renewcommand{\arraystretch}{1.25}%
  \begin{tabularx}{\linewidth}{|L{2.0}|C{0.7}|C{0.7}|C{0.7}|C{0.6}|}
    \hline
    \textbf{Question} & \textbf{Réponse A} & \textbf{Réponse B} & \textbf{Réponse C} & \textbf{Corrigé} \\
    \hline
    \textbf{1.} L'opposé de $-\dfrac{3}{4}$ est :
      & $\dfrac{3}{-4}$ & $-\dfrac{-3}{-4}$ & $\dfrac{-3}{-4}$ & \ans{C} \\
    \hline
    \textbf{2.} La somme de deux nombres relatifs de signes différents...
      & est toujours positive & est parfois positive, parfois négative & est toujours négative & \ans{B} \\
    \hline
    \textbf{3.} Le produit de 163 nombres non nuls dont 47 sont positifs est...
      & parfois positif et parfois négatif & toujours positif & toujours négatif & \ans{B} \\
    \hline
    \textbf{4.} Soient $a>0$, $b<0$ et $c<0$. Alors :
      & $\dfrac{a\times a}{b}>0$ & $\dfrac{b}{-c}>0$ & $\dfrac{-ab}{c}<0$ & \ans{C} \\
    \hline
  \end{tabularx}}
\end{exercisebox}

% ======== Exercice 1 ========
\begin{exercisebox}{\faIcon{edit} Exercice 1 : Opérations de base \points{4}}
  \vspace{0.2cm}
  \textbf{Calculer les expressions suivantes :}
  \begin{multicols}{2}
    \begin{enumerate}[label=\textcolor{primary}{\textbf{\alph*)}},leftmargin=*]
      \item $(-8) + (+5) = $ \ans{$-3$} \points{1}
      \vspace{0.8cm}
      \item $(-3) - (-7) = $ \ans{$4$} \points{1}
      \vspace{0.8cm}
      \item $(+6) \times (-4) = $ \ans{$-24$} \points{1}
      \vspace{0.8cm}
      \item $(-36) \div (+9) = $ \ans{$-4$} \points{1}
    \end{enumerate}
  \end{multicols}
\end{exercisebox}

% ======== Exercice 2 ========
\begin{exercisebox}{\faIcon{layer-group} Exercice 2 : Priorités opératoires \points{8}}
  \vspace{0.3cm}
  \textbf{Calculer en détaillant toutes les étapes :}
  \begin{enumerate}[label=\textcolor{primary}{\textbf{\alph*)}},leftmargin=*]
    \item $A = -4 + 6 \times (-3)$ \points{1}\\
    \anslines[0.50cm]{4}{\begin{align*}
      A &= -4 + 6\times(-3) \\
        &= -4 - 18 \\
        &= -22
    \end{align*}}

    \item $B = 18 - 2 \times (5 - 9)$ \points{1}\\
    \anslines[0.50cm]{4}{\begin{align*}
      B &= 18 - 2\times(5-9) \\
        &= 18 - 2\times(-4) \\
        &= 18 - (-8) \\
        &= 26
    \end{align*}}

    \item $C = \dfrac{-15 + 3}{-4} - 7$ \points{2}\\
    \anslines[0.50cm]{4}{\begin{align*}
      C &= \dfrac{-15 + 3}{-4} - 7 \\
        &= \dfrac{-12}{-4} - 7 \\
        &= 3 - 7 \\
        &= -4
    \end{align*}}

    \item $D = 3 \times [(-4) + 7] - 2 \times (-6 + 3)$ \points{2}\\
    \anslines[0.40cm]{7}{\begin{align*}
      D &= 3\times[(-4)+7] - 2\times(-6+3) \\
        &= 3\times 3 - 2\times(-3) \\
        &= 9 - (-6) \\
        &= 15
    \end{align*}}

    \item $E = \dfrac{(-2)^2 + 4 \times 3}{-5 + 2}$ \points{2}\\
    \anslines[0.50cm]{4}{\begin{align*}
      E &= \dfrac{(-2)^2 + 4\times 3}{-5 + 2} \\
        &= \dfrac{4 + 12}{-3} \\
        &= \dfrac{16}{-3} \\
        &= -\dfrac{16}{3}
    \end{align*}}
  \end{enumerate}
\end{exercisebox}

% ======== Exercice 3 ========
\begin{exercisebox}{\faIcon{lightbulb} Exercice 3 : Résolution de problème \points{4}}
  \vspace{0.01cm}
  \begin{tcolorbox}[colback=blue!5,colframe=primary!30,arc=2mm,boxrule=0.5pt]
    \textbf{Contexte :} Un sous-marin se trouve à une profondeur de $-120$ mètres.
    \begin{itemize}[noitemsep,topsep=5pt]
      \item Il remonte d'abord de 45 mètres
      \item Puis il descend de 30 mètres
      \item Enfin, il remonte de 25 mètres
    \end{itemize}
  \end{tcolorbox}

  \begin{enumerate}[leftmargin=*]
    \item \textbf{Écrire l'expression mathématique} qui permet de calculer la profondeur finale : \points{1.5}\\
    \anslines[0.40cm]{3}{$-120 + 45 - 30 + 25$}

    \item \textbf{Calculer} cette profondeur finale : \points{1.5}\\
    \anslines[0.40cm]{4}{\begin{align*}
      -120 + 45 &= -75 \\
      -75 - 30  &= -105 \\
      -105 + 25 &= -80
    \end{align*}}

    \item \textbf{Question :} Le sous-marin veut atteindre $-50$ mètres. Doit-il monter ou descendre ? De combien de mètres ? \points{1}\\
    \anslines[0.40cm]{3}{Pour passer de $-80$ m à $-50$ m, il faut monter de $30$ m.}
  \end{enumerate}
\end{exercisebox}

\vspace{0.2cm}

% ======== Bonus ========
\begin{tcolorbox}[
  colback=accent!10,
  colframe=accent,
  fonttitle=\bfseries,
  title={\faIcon{star} BONUS : Défi \badge{accent}{+2pts}},
  arc=3mm,
  boxrule=1.5pt,
  enhanced,
  width=\linewidth,
  grow to left by=3mm,
  grow to right by=3mm
]
  Soit l'expression : $G = \dfrac{2a - b^2}{c + d}$\\
  Avec : $a = -4$ ; $b = -3$ ; $c = -5$ ; $d = 2$\\
  Calculer $G$ en remplaçant et en détaillant toutes les étapes.\\
  \anslines[0.40cm]{7}{\begin{align*}
    G &= \dfrac{2\times(-4) - (-3)^2}{-5 + 2} \\
      &= \dfrac{-8 - 9}{-3} \\
      &= \dfrac{-17}{-3} \\
      &= \dfrac{17}{3}
  \end{align*}}
\end{tcolorbox}

% ======== Pied de page ========
\vspace{0.2cm}
\begin{center}
  \begin{tikzpicture}
    \node[draw=secondary,rounded corners=3mm,inner sep=8pt,fill=light!30] {
      \begin{tabular}{lccccc}
        \textbf{Auto-évaluation :} &
        \faIcon{frown} Difficile &
        \faIcon{meh} Moyen &
        \faIcon{smile} Facile &
        \phantom{xx} &
        \textbf{Temps utilisé :} \underline{\phantom{xxxxx}} min
      \end{tabular}
    };
  \end{tikzpicture}
\end{center}

